\uuid{Ik2A}
\titre{Intégration/DL usuels/Dl en l'infini}
\auteur{hanani ; mzari}
\organisation{exo7}
\licence{CC BY-NC-SA 4.0}
\datecreate{2020-01-20}
\niveau{L1}
\module{Analyse}
\chapitre{Développement limité}
\sousChapitre{Calculs}
\difficulte{3}
\theme{Opérations sur les DL, Intégration, DL usuels, Dl en l'infini}
\exo7qcmid{lille-2.developpements-limites.028}

\contenu{
\texte{
Soit $a\in \Rr$ et $f$ une fonction. On dit que $f(x)$ est un $o\left((x-a)^n\right)$ au voisinage de $a$ et on écrit $f(x)=o((x-a)^n)$ si $\displaystyle \lim _{x\to a}\frac{f(x)}{(x-a)^n}=0$. On note par $DL_n(a)f(x)$ le développement limité de $f$ à l'ordre $n$ en $a$. On note aussi $DL_n(a^+)f(x)$ (resp. $DL_n(a^-)f(x)$) le développement limité de $f$ à l'ordre $n$ à droite en $a$ (resp. à gauche en $a$).
}

\begin{question}
\qclassification{125.02}
\qtags{motcle=Intégration/DL usuels/Dl en l'infini}
Soit $\displaystyle f(x)=\arctan\left(\frac{x}{x+1}\right)$. On considère la fonction $g$ définie par $\displaystyle g(t)=f\left(\frac{1}{t}\right)$. Parmi les affirmations suivantes, cocher celles qui sont vraies :
\begin{answers}
\good{Pour $0<t$, $\displaystyle g(t)=\frac{\pi}{2}-\arctan (1+t)$ et $\displaystyle DL_1(0)g'(t)=-\frac{1}{2}+\frac{t}{2}+o(t)$.}
\bad{$\displaystyle DL_2(0^+)g(t)=-\frac{t}{2}+\frac{t^2}{4}+o(t^2)$.}
\bad{$\displaystyle DL_2(+\infty)f(x)=-\frac{1}{2x}+\frac{1}{4x^2}+o\left(\frac{1}{x^2}\right)$.}
\good{$\displaystyle DL_2(+\infty)f(x)=\frac{\pi}{4}-\frac{1}{2x}+\frac{1}{4x^2}+o\left(\frac{1}{x^2}\right)$.}
\end{answers}
\begin{explanations}
Pour $0<t$, on a : $\displaystyle g(t)=\arctan\left(\frac{1}{1+t}\right)=\frac{\pi}{2}-\arctan (1+t)$ et
$$g'(t)=\frac{-1}{1+(1+t)^2}=\frac{-1}{2}\times \frac{1}{1+(t+t^2/2)}=-\frac{1}{2}+\frac{t}{2}+o(t).$$
Par intégration, et puisque $\displaystyle \arctan 1=\frac{\pi}{4}$, on obtient :
$$DL_2(0^+)g(t)=\frac{\pi}{4}-\frac{t}{2}+\frac{t^2}{4}+o(t^2)\Rightarrow DL_2(+\infty)f(x)=\frac{\pi}{4}-\frac{1}{2x}+\frac{1}{4x^2}+o\left(\frac{1}{x^2}\right).$$
\end{explanations}
\end{question}
}
